% Source: https://artofproblemsolving.com/community/c3h1289461_problem_with_coprime_integers
% Stable identity: aops-topic-1289461-q01
If a prime $p$ divides two of the five numbers, it divides their difference, which is $1,2,3$ or $4$. So $p=2$ or $p=3$.

Among five consecutive integers, at most one residue class modulo $6$ is missing, so at least one of the classes $1$ and $5$ modulo $6$ occurs. Let $m$ be a number in one of those classes; then $m$ is divisible by neither $2$ nor $3$.

Any common prime factor of $m$ and another of the numbers would have to be $2$ or $3$, and $m$ has neither. Hence $m$ is relatively prime to each of the other four. $\blacksquare$
